% V.2d and V.3: the results, their critical discussion and the conclusions are
% what counts towards the page minimum, together with the introduction.
%
% The tables and figures of this chapter come from report/tables/generated/ and
% report/figures/generated/, which `musicsim export` refreshes. They are
% versioned on purpose, so that the report compiles without re-running any
% experiment.
%
% Section order follows the question-to-experiment traceability of this work:
% the siamese-network section, its principal focus, comes first, then the rest
% in the order the questions are introduced in the introduction chapter.

\chapter{Results}
\label{cap:results}

\section{Siamese networks}
\label{sec:results-siamese}

\todo{SQ5, the principal focus of the work. Experiments e05, e08, e10. Report,
in order: (1) retrieval and structure for each siamese input
(\texttt{siamese\_mel}, \texttt{siamese\_emb}, \texttt{siamese\_desc}) and for
\texttt{triplet\_mel}, with parameter counts; (2) the gain of each siamese
network over its own raw input (\texttt{siamese\_emb} against \texttt{mert},
\texttt{siamese\_desc} against \texttt{acoustic}) as the headline SQ5 result;
(3) siamese networks against the classical block (\texttt{mfcc},
\texttt{acoustic}, \texttt{gauss\_mfcc}+KL, best \texttt{posterior\_<clf>}) and
against the CNN; (4) the effect of the similarity measure (e10) on the
classical representations and on the best siamese model; (5) classifier
accuracy and macro F1 of the classical classifiers on \texttt{fma\_test}. Every
comparison follows the Holm-corrected family declared in the evaluation
protocol.}

\section{Absolute quality}
\label{sec:results-absolute}

\todo{SQ1. Experiments e01, e02, e04, e07. MFCC against the random
control on FMA, for the three exclusion variants and both query protocols,
with intervals by query and by artist, and the per-genre breakdown (e01);
paired differences under PCA against the untouched embedding (e02); agreement
with the MTT triplets against the reference points of the evaluation protocol,
and the faceted relevance results on MTT (e04); robustness of the neighbour sets
under perturbation (e07).}

\subsection{Structure of the graph}
\label{sec:results-structure}

\emph{Proxy-based (genre) evidence only. Provisional: figures from the e03 run
of 6 November 2026; they are regenerated into \texttt{report/tables/generated}
before submission.} \todo{Replace the numbers in this subsection by the generated
table, and add the three figures of the run (k-occurrence distribution, community
sizes, community $\times$ genre matrix).}

On \texttt{fma\_test} (795 nodes after collapsing 5 near-duplicates, $k=10$)
the MFCC graph is clearly more structured than the random control, on the three
headline metrics (Table~\ref{tab:structure-test}). Its k-occurrence is
right-skewed (skewness 2.37, interval over artist subsamples [1.96, 2.62]) where
that of a random embedding is almost symmetric (0.18, [0.05, 0.33]); label
homophily is 0.35 [0.33, 0.39] against 0.12 [0.12, 0.15], and the NMI between
Louvain communities and genre is 0.19 [0.15, 0.22] against 0.03 [0.03, 0.04].
The intervals do not overlap. Both graphs are one connected component. As
expected, the neighbour-vote accuracy follows homophily (0.24 against 0.07),
and is a description of the graph, not a classifier score
(Section~\ref{sec:structure-protocol}).

\begin{table}[t]
\centering
\caption{Headline structure metrics on \texttt{fma\_test}, $k=10$, 795 nodes.
Interval: 2.5th--97.5th percentile over 200 paired subsamples of 80\,\% of the
artists, graph rebuilt in each; a sensitivity interval, not a confidence
interval. Proxy-based family. Provisional.}
\label{tab:structure-test}
\begin{tabular}{lccc}
\hline
 & Skewness of $N_{10}$ & Label homophily & NMI Louvain--genre \\
\hline
\texttt{mfcc}      & 2.37 [1.96, 2.62] & 0.35 [0.33, 0.39] & 0.19 [0.15, 0.22] \\
\texttt{random\_s0} & 0.18 [0.05, 0.33] & 0.12 [0.12, 0.15] & 0.03 [0.03, 0.04] \\
\hline
\end{tabular}
\end{table}

Three points qualify these numbers.
\begin{itemize}
  \item \textbf{Homophily is not a second retrieval result.} It equals the P@10
        of the \texttt{raw} variant, and the random control's value (0.12) is
        the chance level for eight genres.
  \item \textbf{The whole corpus gives the same picture at a larger scale.} On
        \texttt{fma\_all} (7978 nodes, point estimates only) the skewness is 4.70
        for \texttt{mfcc} against about 0.3 for the random draws, homophily
        0.327 against 0.125, and the giant component holds all the nodes. The NMI is
        0.133 for \texttt{mfcc}, below the range of 0.15 to 0.30 the plan
        expected; the Louvain resolution was left at 1 and not tuned to enter
        that range. \todo{Say whether this is read as a weaker alignment between
        timbre communities and genre at this scale, after the per-community genre
        matrix has been inspected.}
  \item \textbf{Intervals are not centred on the point estimate.} The subsamples
        have 80\,\% of the artists and therefore fewer nodes, which changes the
        k-occurrence and the communities; the estimate on the full gallery can
        lie at an edge of its interval or, for the comparison between graphs
        below, outside it.
\end{itemize}

Reducing MFCC with PCA changes the graph progressively. The neighbour Jaccard@10
between the MFCC graph and its PCA reduction is 0.94 at 64 dimensions, 0.80 at 48,
0.64 at 32 and 0.40 at 16, while it is 0.007 to 0.008 between MFCC and any of the
three random draws, and 0.008 among the random draws themselves. The latter is
the floor of a graph without structure; whether the differences between
\emph{trained} representations exceed the noise between seeds of the same model
is a question for the criterion of Section~\ref{sec:pq-criterion} and cannot be
answered with the representations available so far. \todo{PCA deltas of the
headline metrics against the untouched embedding, with their paired intervals,
from \texttt{structure\_deltas.csv}; \texttt{mert} rows when the model is
extracted on both galleries.}

\section{Model comparison}
\label{sec:results-comparison}

\todo{SQ2. Experiment e08. The master table: every representation, every
seed, against every headline metric, with Holm-corrected paired differences and
a forest plot, on \texttt{fma\_test}, \texttt{mtt\_all} and
\texttt{mtt\_triplets}.}

\section{Representation gap and graph gap}
\label{sec:results-link}

\todo{SQ3. Experiment e09. The noise floor: Jaccard@10 and CKA between seeds of
the same model against the same quantities between different models. The
dose-response scatter of $1-\text{CKA}$ against $1-\text{Jaccard@10}$ over the
$\approx$45 representation points of the evaluation protocol, with the Mantel
test. The descriptive
comparison of probe score (e05) against graph quality and structure.}

\section{Fuzzy modelling}
\label{sec:results-fuzzy}

\todo{SQ6. Experiment e11. Crisp against fuzzy graphs for \texttt{mfcc}, the
best siamese model and \texttt{mert}: retrieval ordered by membership degree,
the curve of metrics across $\alpha$-cuts, fuzzy ARI against the multi-label
ground truth compared with Louvain's ARI, the model-versus-human ambiguity
correlation (F4), and, if F2 is in scope, the fitted rule base.}

\section{Proxy versus perception}
\label{sec:results-proxy}

\todo{SQ4. Experiments e04, e08. Kendall $\tau$ between the genre-based and the
human-judgement-based rankings of the representations, and whether the SQ2
conclusions drawn on FMA survive on MTT.}

\section{Sensitivity}
\label{sec:results-sensitivity}

\todo{Robustness questions. Experiments e06, e07. kNN against mutual kNN
against the $\varepsilon$ threshold, with and without mutual proximity (e06);
neighbour stability under each perturbation (e07); state whether the sign of
the headline deltas changes under any crisp graph variant.}
